\documentclass[11pt]{amsart}
\usepackage[margin=1.1in]{geometry}
\usepackage{amsmath,amssymb,amsthm}
\usepackage[colorlinks=true,linkcolor=blue!60!black,citecolor=blue!60!black,urlcolor=blue!60!black]{hyperref}
\usepackage{enumitem}
\usepackage{tikz-cd}

\newtheorem{theorem}{Theorem}[section]
\newtheorem{proposition}[theorem]{Proposition}
\newtheorem{lemma}[theorem]{Lemma}
\theoremstyle{remark}
\newtheorem*{remark}{Remark}
\newtheorem*{caveat}{Caveat}

\newcommand{\C}{\mathbb{C}}
\newcommand{\Z}{\mathbb{Z}}
\newcommand{\N}{\mathbb{N}}
\newcommand{\PP}{\mathbb{P}}
\newcommand{\OO}{\mathcal{O}}
\newcommand{\KK}{\mathcal{K}}
\newcommand{\Gr}{\mathrm{Gr}}
\newcommand{\GL}{\mathrm{GL}}
\newcommand{\End}{\mathrm{End}}
\newcommand{\Hom}{\mathrm{Hom}}
\newcommand{\Spec}{\mathrm{Spec}}
\newcommand{\Hilb}{\mathrm{Hilb}}
\newcommand{\eu}{\mathrm{eu}}
\newcommand{\bv}{\mathbf{v}}
\newcommand{\bw}{\mathbf{w}}
\newcommand{\bb}{\mathbf{b}}
\newcommand{\bd}{\mathbf{d}}
\newcommand{\Ttop}{T_{\mathrm{top}}}
\newcommand{\Hc}{H^*_F(X)}

\title{A reader's guide to\\ ``The equivariant cohomological Hikita conjecture for arbitrary quivers''}
\author{Notes prepared for the SLMath QFT program}
\date{October 2026}

\begin{document}
\maketitle

\begin{caveat}
These notes summarize a preprint released by OpenAI on 24 September 2026
\cite{OAI-Hikita} as part of a collection of manuscripts written by an
unreleased internal model.  The preprint has not been refereed and has no
Lean formalization in the release; the release notes themselves say
that unformalized results ``could have issues.''  These notes were drafted
by Claude (Anthropic) at Ben Webster's request from the \LaTeX\ source,
and reorganize the argument for readability.  They do \emph{not} verify
the proofs: every ``proof'' below is a sketch of what the preprint
claims, and the words ``the paper shows'' should be read as ``the paper
asserts, with an argument I have summarized but not checked.''
Section~\ref{sec:assessment} records which inputs are standard and
which steps are new.
\end{caveat}

\section{The statement}

\subsection{Data}
Fix a finite quiver $Q=(Q_0,Q_1)$, loops and multiple arrows allowed, with
gauge dimension vector $\bv$ and framing vector $\bw$.  Put
\[
 G=\prod_{i\in Q_0}\GL(V_i),\qquad
 N=\bigoplus_{a:i\to j}\Hom(V_i,V_j)\ \oplus\ \bigoplus_i\Hom(V_i,W_i),
\]
with $\dim V_i=v_i$, $\dim W_i=w_i$.  Let $F$ be any torus acting on $N$
and commuting with $G$ (a ``flavor torus''), and let $\mu:T^*N\to\mathfrak g^*$
be the moment map.

Fix a character $\nu$ of $G$.  The paper calls $\nu$ \emph{regular} if, on
$\mu^{-1}(0)$, every $\nu$-semistable point is stable and $G$ acts freely
on the semistable locus.  Then
\[
 X=\mu^{-1}(0)^{\nu\text{-st}}/G
\]
is a smooth Nakajima quiver variety, possibly empty.

\subsection{The two sides}
\emph{Higgs side.}  $H^*_F(X;\C)$, the $F$-equivariant cohomology.

\emph{Coulomb side.}  Let $\OO=\C[[t]]$, $\KK=\C((t))$, and let
$\mathcal R=\{[g,s]\in G(\KK)\times_{G(\OO)}N(\OO): gs\in N(\OO)\}$ be the
Braverman--Finkelberg--Nakajima space of triples.  Its
$G(\OO)\rtimes F$-equivariant Borel--Moore homology
\[
 A_F=H_*^{G(\OO)\times F}(\mathcal R)
\]
is a commutative algebra under convolution (there is no loop-rotation
parameter, so this is the \emph{commutative} Coulomb branch, deformed by
the flavor parameters).  Its spectrum $Y_{\mathfrak f}=\Spec A_F$ is the
flavor-deformed Coulomb branch.  The algebra $A_F$ is graded by
$\pi_1(G)=\Z^{Q_0}$ (``topological charge'' or ``monopole charge''), which
is the same as an action of the torus $\Ttop=\Hom(\pi_1(G),\C^\times)$.
The character $\nu$ of $G$ is a cocharacter of $\Ttop$, and the
\emph{scheme-theoretic} fixed locus of this cocharacter has coordinate ring
\[
 \C[Y_{\mathfrak f}^\nu]=A_F\big/\big\langle (A_F)_m : m\neq 0\big\rangle,
\]
where $(A_F)_m$ is the subspace on which $\nu$ acts with weight $m$, i.e.\
the sum of charge spaces $A_F^{(d)}$ with $\nu\cdot d=m$.  
taken: nilpotents are kept.

\subsection{The common coefficient ring}
Both sides are algebras over
\[
 R=H^*_{G\times F}(\mathrm{pt})=\C[\mathfrak t\oplus\mathfrak f]^W,
\]
the ring of symmetric functions in the Chern roots $x_{ia}$ of the gauge
bundles together with the flavor parameters.  On the Higgs side $R$ acts by
the Kirwan map $\phi_H:R\to \Hc$ (Chern classes of tautological
bundles).  On the Coulomb side $R=H^*_{G(\OO)\times F}(\mathrm{pt})$ is the
charge-zero part of $A_F$, so it maps to $\C[Y_{\mathfrak f}^\nu]$ by
$\phi_C$.  The content of the conjecture is that the \emph{relations}
agree.

\begin{theorem}[{\cite[Theorem~1.1]{OAI-Hikita}}]\label{thm:main}
For every finite quiver, every $\bv,\bw$, every commuting flavor torus $F$
and every regular character $\nu$, both $\phi_H$ and $\phi_C$ are
surjective with the same kernel.  Hence there is a canonical isomorphism of
graded $R$-algebras
\[
 \Hc\ \cong\ \C[Y_{\mathfrak f}^\nu].
\]
The grading on the left is cohomological degree, on the right the BFN
homological degree.  The empty case ($X=\varnothing$, both rings zero) is
included.
\end{theorem}

This is the conjecture as formulated by Dumanski--Krylov
\cite[Conj.~2.1, eq.~(2.3)]{DK25}, building on Hikita \cite{Hikita17},
Nakajima, and Kamnitzer--Tingley--Webster--Weekes--Yacobi.  Previously known cases:
the Jordan quiver (Gieseker varieties, Krylov--Shlykov \cite{KS25}) and
finite type $A,D,E_6$, with restrictions on the framing in $E_7,E_8$
(Dumanski--Krylov \cite{DK25}).  The theorem does \emph{not} address the
quantized version, nor gauge groups other than products of $\GL$'s.


\section{Architecture of the proof}

Introduce the intermediate ring
\[
 B_F=\C[Y_{\mathfrak f}^{\Ttop}]=A_F\big/\big\langle A_F^{(d)} : d\neq 0\big\rangle,
\]
the coordinate ring of the fixed scheme of the \emph{full} topological
torus (equivalently: kill every monopole operator of nonzero charge).
There are natural maps
\[
 \C[Y_{\mathfrak f}^\nu]\ \twoheadrightarrow\ B_F,
 \qquad\qquad
 R\ \twoheadrightarrow\ B_F,\quad R\ \twoheadrightarrow\ \Hc .
\]
The proof splits into two independent halves.

\medskip
\noindent\textbf{Part I} ($\S\S3$--$7$ of the paper): \ $B_F\cong\Hc$, i.e.\
$\ker(R\to B_F)=\ker(R\to\Hc)$.

\noindent\textbf{Part II} ($\S8$): \ $\C[Y_{\mathfrak f}^\nu]\to B_F$ is an
isomorphism, i.e.\ the ideal generated by nonzero-$\nu$-weight monopoles
already contains every nonzero-charge monopole.

\medskip
Part I is itself a three-stage argument: prove one inclusion of ideals in
general; prove equality in a special situation (loops at every vertex,
determinant stability) by a size count; deform back.

\section{Part I, stage 0: generators on both sides}

Both $B_F$ and $\Hc$ are quotients of $R$:
\begin{itemize}[leftmargin=2em]
\item \emph{Higgs:} Kirwan surjectivity for quiver varieties
(McGerty--Nevins \cite{MN18}), including the equivariant version.  This also
gives $H^{\mathrm{odd}}(X)=0$.
\item \emph{Coulomb:} BFN abelianization embeds $A_F$ into
$\big(\mathrm{Frac}(P_F)[X_*(T)]\big)^W$ where $P_F=\C[x_{ia}]\otimes\C[\mathfrak f]$,
and Weekes's generation theorem \cite{Weekes19} says $A_F$ is generated over $R$ by
the \emph{dressed minuscule monopoles}: for a coweight $\lambda$ all of whose
entries are in $\{0,1\}$ (or all in $\{0,-1\}$) and a $W_\lambda$-invariant
polynomial $p$,
\begin{equation}\label{eq:minuscule}
 M_{\lambda,p}=\sum_{\lambda'\in W\lambda}
 e^{\lambda'}\,p_{\lambda'}\,
 \frac{\prod_{\beta\in\mathrm{wt}(N)}\beta^{(-\langle\beta,\lambda'\rangle)_+}}
      {\eu_T\big(T_{t^{\lambda'}}\Gr_G^{\lambda}\big)} .
\end{equation}
Consequently the ideal $\langle A_F^{(d)}:d\neq0\rangle$ is generated by the
nontrivial minuscule monopoles, and any relation $u\in\ker(R\to B_F)$ can be
written as a finite sum of \emph{neutral words}: products of minuscule
monopoles of total charge zero.
\end{itemize}

\section{Part I, stage 1: the stable-pair test (one inclusion)}

\begin{proposition}[{\cite[Prop.~4.1]{OAI-Hikita}}]
For all data in Theorem~\ref{thm:main}, $\ker(R\to B_F)\subseteq\ker(R\to\Hc)$;
so the coefficient map induces a graded surjection $B_F\twoheadrightarrow\Hc$.
\end{proposition}

\noindent\emph{Idea.}  Take a Coulomb relation $u\in\ker(R\to B_F)$.  By
stage~0, $u$ lies in the ideal of $A_F$ generated by the nontrivial minuscule
monopoles, and since $u$ has charge zero we may discard all other charge
components, so
\[
 u=\sum_w M_w,\qquad
 M_w=M_{\lambda_{w,1},p_{w,1}}\cdots M_{\lambda_{w,k_w},p_{w,k_w}},
\]
where every letter is a nontrivial (hence nonzero-charge) minuscule
monopole and the charges (images in $\pi_1(G)$) of the letters in each word
sum to zero.  There are no separate coefficients: for $r\in R$, which is
$W$-invariant, \eqref{eq:minuscule} gives $r\,M_{\lambda,p}=M_{\lambda,rp}$,
so coefficients are absorbed into the dressings.
One wants to show $u$ restricts to zero on $D=\mu^{-1}(0)^{\nu\text{-st}}$.
The paper does this by building a finite-dimensional ``test'' that sees both
the identity lattice (where $u$ acts by cup product) and the first step of
each word (which it kills).

\begin{enumerate}[leftmargin=2em]
\item \emph{Finite correspondences} (\S3).  Truncate $N(\KK)$ to
$L=t^{-l}N(\OO)/t^lN(\OO)$ and, for a bounded family $B$ of lattices
$b=gG(\OO)$, let $S_B\to L$ be the family of truncated lattices
$E_b=gN(\OO)/t^lN(\OO)$.  A word in minuscule monopoles becomes a chain of
incidence correspondences between such families, and the paper proves
(by a localization argument that is faithful because the relevant
Borel--Moore homology is free over $H^*_{T\times F}(\mathrm{pt})$) that the
relation $u=\sum_w M_w$ in $A_F$ becomes an identity
$u\cdot\rho_Z=\sum_w\Psi_w$ of morphisms of constructible sheaves on $L$,
where $\rho_Z$ is restriction to the identity lattice and each $\Psi_w$
factors through the family of the \emph{first} lattice of the word $w$.
\item \emph{The test} (\S4).  On $\widetilde L=L\times D\times P_-$ define
the function
\[
 W(n,s,p,y)=\mathrm{Res}\big(\langle p+y,n\rangle-\langle y,s\rangle\big),
 \qquad (s,p)\in D,
\]
and test a sheaf $\mathcal F$ on $L$ by taking $G\times F$-equivariant
cohomology of the vanishing cycles $\Phi_W(\mathrm{pr}^*\mathcal F)$.  The
critical locus of $W$ restricted to a lattice family is exactly the locus
where the stable pair $(s,p)$ extends across the lattice modification to a
pair of sections of the associated bundle on $\PP^1$.
\item \emph{Stability kills the first step.}  If $\mathcal P$ is a
$G$-bundle on $\PP^1$ with a section of $\mathcal P\times_G T^*N$ whose
value at $\infty$ is $\nu$-stable with trivial stabilizer, then
$\deg\mathcal L_\nu\ge0$, with equality only if $\mathcal P$ is trivial
and the section is constant.  (Proof: a $\nu^m$-semi-invariant nonzero at
$\infty$ gives a nonzero section of $\mathcal L_\nu^{m}$; in degree zero
the induced map $\PP^1\to T^*N/\!/_\nu G$ is constant because the GIT
polarization is relatively ample over the affine quotient.)  Since each
word is neutral, the $\nu$-degrees of its steps sum to zero, so (by
commutativity) one can put a step of nonpositive $\nu$-degree and nonzero
charge first.  Its critical locus is empty, so its test vanishes.
\item \emph{The identity lattice survives.}  On the neutral family the test
is computed in a quadratic normal form: after a change of variables the
function becomes $\langle y',a'\rangle$ (this uses $\mu(s,p)=0$ to show the
critical value is zero), so vanishing cycles are a rank-one local system
shifted by $2r$, and the test of $\rho_Z$ becomes ordinary restriction
$H^*_{G\times F}(U\times D)\to H^*_{G\times F}(D)$, which is split surjective.
\end{enumerate}
Combining: $u|_D\smile j^*=0$ with $j^*$ split, so $u|_D=0$ in
$H^*_{G\times F}(D)=\Hc$.

\section{Part I, stage 2: add loops and count}

Now suppose temporarily that $Q$ has a loop at every vertex with $v_i>0$,
there is no flavor, and $\nu=\nu_{\det}=\prod_i\det_{V_i}$ (which is always
regular).  Write $B_0$ for the no-flavor version of $B_F$.

\subsection{The Laurent model}  The key observation is that one loop at a
vertex contributes to the numerator of \eqref{eq:minuscule} the factor
$\prod_{a\in I,b\notin I}(x_{ib}-x_{ia})$, which is exactly the Euler class
of the tangent space of the Grassmannian $\Gr^{\lambda}$ at the fixed point
$t^{\lambda'}$.  So with a loop at every vertex, the denominators in
\eqref{eq:minuscule} cancel and $A_0$ embeds in a \emph{polynomial} Laurent
ring:
\begin{proposition}[{\cite[Prop.~5.1]{OAI-Hikita}}]
Let $\mathcal B'$ be the multiset of weights of $N$ with one loop per vertex
removed, and set $r_\lambda=e^\lambda\prod_{\beta\in\mathcal B'}\beta^{(-\langle\beta,\lambda\rangle)_+}$.
Then
\[
 A_0=\Big(\bigoplus_{\lambda\in\Z^{\Omega}}P\,r_\lambda\Big)^{W},
 \qquad P=\C[x_{ia}].
\]
\end{proposition}
Both inclusions need proof; the nontrivial one ($\supseteq$) is an
induction on $\|\lambda\|^2$ using a rearrangement-inequality argument to
show that each $r_\lambda$ with arbitrary stabilizer-invariant dressing is
a product of two shorter ones plus lower terms.

\subsection{Filtration by multipartitions and cut ideals}
A multipartition $m=(m_{i,s})$ of $\bv$ (i.e.\ $\sum_s s\,m_{i,s}=v_i$)
determines a charge-zero coweight $\lambda_m$ with entries $s-1$ on
$m_{i,s}$ slots of color $i$.  The paper filters $B_0$ by the number $|l|$
of nonzero parts and shows that the associated graded piece for $m$ is a
quotient of
\[
 (P_m/J_m)^{W_m},\qquad P_m=\C[x_d: d\text{ a nonzero part}],
\]
where $J_m$ is the \emph{cut ideal} generated by
\[
 q_S=\prod_{d\in S}x_d^{w_{i(d)}}\prod_{d\in S,\ e\notin S}(x_d-x_e)^{b_{de}},
 \qquad \varnothing\neq S\subseteq\{\text{parts}\},
\]
with $b_{de}=c_{i(d)i(e)}\min(s(d),s(e))$ and $c_{ij}$ the number of arrows between $i,j$ in either direction (after removing the chosen loops).
The relations $q_S$ come from the vanishing in $B_0$ of products
$M_\alpha M_\beta$ with $\beta$ of nonzero positive charge: the excess
weight factor in $r_\alpha r_\beta/r_{\lambda_m}$ is precisely $q_S$.
These cut ideals are the cocircuit ideals familiar from toric hyperk\"ahler
cohomology (Hausel--Sturmfels).  The upshot is a coefficientwise bound
\begin{equation}\label{eq:bound}
 \Hilb(B_0;q)\ \le\ \sum_{m}q^{\delta_m}\,\Hilb\big((P_m/J_m)^{W_m};q\big),
\end{equation}
with explicit shifts $\delta_m$.

\subsection{Matching Hausel's formula}
The paper computes $\Hilb((P_m/J_m)^{W_m})$ by realizing $P_m/J_m$ as the
$T_m$-equivariant cohomology of an open ``all slots reachable from the
framing'' locus in a representation space, stratified by the set of
unreachable slots.  This gives a recurrence whose generating function is a
ratio of two $q$-series, and that ratio is exactly the right-hand side of
Hausel's arithmetic-Fourier-transform formula \cite{Hausel06} for the
Poincar\'e polynomial of $\mu^{-1}(\mathrm{id})/\!/G$.  A hyperk\"ahler
rotation identifies that affine quotient with $X_{\det}$ topologically.
Hence the right side of \eqref{eq:bound} equals $P_{X_{\det}}(q)$.
Together with the surjection of stage 1:

\begin{proposition}[{\cite[Cor.~6.6]{OAI-Hikita}}]
With a loop at every vertex and $\nu=\nu_{\det}$, the coefficient map is an
isomorphism $B_0\cong H^*(X_{\det})$.
\end{proposition}

\section{Part I, stage 3: remove the loops, then vary everything}

\subsection{Loop removal by a mass deformation}
For an arbitrary quiver $Q$, let $Q^+$ add one loop at each vertex, and let
an auxiliary $H=\C^\times$ act with weight $1$ on the new loops only, with
equivariant parameter $z$.  The previous stage plus a flavor-lift lemma
(see below) gives $B^+_H\cong H^*_H(X^+)$ as $R_0[z]$-algebras.  Now localize
at the maximal ideal of $(x,z)=(0,1)$, i.e.\ near zero gauge parameters and
mass one.
\begin{itemize}[leftmargin=2em]
\item \emph{Coulomb side:} the extra loop contributes to each minuscule
monopole the factor $E_\lambda=\prod(z+x_{ib}-x_{ia})^{(\lambda_{ia}-\lambda_{ib})_+}$,
which equals $1$ at $(0,1)$, hence is a unit after localization.  So the
localized Coulomb algebras of $Q$ and $Q^+$ agree, compatibly with the
charge grading.
\item \emph{Higgs side:} Brion's localization theorem on the stable
prequotient $D^+$ (not on the quotient, where $H$-fixed points need
compensating gauge transformations) shows restriction to $(D^+)^H=D$ is an
isomorphism after inverting characters nontrivial on $H$.  Stability
restricts correctly under the closed embedding $T^*N\subset T^*N^+$.
\end{itemize}
This gives equality of the two ideals in $R_0$ after localizing at the
graded origin.  Both ideals are homogeneous, and a homogeneous ideal is
determined by its localization at the origin (if $sh\in I$ with
$s(0)\neq0$, take the component of degree $\deg h$).  So the ideals are
equal globally.  This is Prop.~7.3: the point is that one does not need
to know that a ``generic'' isomorphism specializes.

\subsection{Other stabilities}  All regular $\nu$ give quotients with the
same Betti numbers (Lemma 7.1): a hyperk\"ahler rotation identifies
$X_\nu$ with the affine quotient $\mu^{-1}(c\nu)/\!/G$, and the compactly
supported Betti numbers of the fibers of $\{\mu=\zeta\}/\!/G\to\mathfrak c$
are constant on a dense open set that every real ball meets.  So for any
regular $\nu$, the surjection $B_0\twoheadrightarrow H^*(X_\nu)$ of stage 1
is between spaces of the same finite dimension.

\subsection{Flavor lift}  If $B_0\cong H^*(X)$ then $B_F\cong\Hc$ for any
$F$: $\Hc$ is free over $\C[\mathfrak f]$ (Leray--Hirsch, using
$H^{\mathrm{odd}}=0$), $B_F$ base-changes to $B_0$ (an even-strata
argument on the space of triples), and a graded Nakayama argument kills
the kernel.

\medskip
This completes Part I: $B_F\cong\Hc$ for all data.

\section{Part II: cocharacter-fixed versus torus-fixed}

\begin{proposition}[{\cite[Prop.~8.4]{OAI-Hikita}}]
For regular $\nu$, the quotient map $\C[Y_{\mathfrak f}^\nu]\to B_F$ is an
isomorphism.
\end{proposition}
Equivalently: every nontrivial minuscule monopole already lies in the ideal
generated by monopoles of nonzero $\nu$-weight.  If $\nu\cdot\bb\neq0$ for all
$0<\bb\le\bv$ this is immediate; the content is the hyperplane case.

\noindent\emph{Idea: complementary monopoles.}  Adjoin a framing vertex
$\infty$ of dimension $1$ with $w_i$ arrows $i\to\infty$, making an
\emph{unframed} quiver $\widehat Q$ with dimension vector $\bd=(\bv,1)$ and
the same representation space.  The extra $\C^\times$ is central and acts
trivially, so
\[
 A_{\widehat Q}\cong A_F\otimes\C[\xi,z^{\pm1}],
\]
where $z$ is the invertible monopole of charge $\bd$ (``translate every
coweight entry by one'').  Extend $\nu$ by $\widehat\nu_\infty=-\sum\nu_iv_i$
so that $\widehat\nu\cdot\bd=0$; then $z$ has weight zero and the cocharacter
quotients are compatible.  Translation by $z$ gives the identity
\[
 M^+_{\bb,p}=z\,M^-_{\bd-\bb,p}
\]
between a positive minuscule monopole of charge $\bb=(\bb',0)$ and the
\emph{complementary} negative one.  So if $M^+_{\bb,p}$ survived in the
$\widehat\nu$-fixed quotient, so would $M^-_{\bd-\bb,p}$.

Now use the polynomial-form cohomological Hall algebra of the tripled
quiver (Jindal--Negu\c t), together with Tsymbaliuk's difference-operator
realization of the shuffle algebra, to prove (Lemma 8.2): every minuscule
monopole of charge $\bb$ is an $R$-combination of words whose factors have
charges $\gamma_h\in\Delta_{\widehat Q}$ (the support of Kac's polynomials
$A_{\widehat Q,\gamma}(q)$, i.e.\ dimension vectors admitting absolutely
indecomposable representations) summing to $\bb$.  A word survives the
cocharacter quotient only if each factor has $\widehat\nu$-weight zero.  So
survival of both complementary monopoles decomposes $\bd=\sum_h\gamma_h$
into at least two Kac-support vectors each orthogonal to $\widehat\nu$.

Finally (Lemma 8.3, essentially Crawley-Boevey's lifting lemma): for
$\gamma\in\Delta_{\widehat Q}$ with $\widehat\nu\cdot\gamma=0$, the
hyperk\"ahler moment fiber $\mu_\C^{-1}(0)\cap\mu_{\mathbb R}^{-1}(\tau_{\widehat\nu})$
is nonempty.  The direct sum of such points for the blocks $\gamma_h$ is
$\widehat\nu$-semistable with a stabilizer containing $k\ge2$ independent
scalar circles; after removing the ineffective diagonal one, the stabilizer
is still positive-dimensional.  This contradicts regularity of $\nu$
(including when the stable locus is empty: it would produce a semistable
point where there is none).  Hence no nontrivial minuscule monopole
survives, and the two ideals coincide.

\section{Assessment}\label{sec:assessment}

\subsection*{Inputs from the literature}
BFN abelianization and minuscule generation (BFN, Weekes); Kirwan
surjectivity for quiver varieties (McGerty--Nevins); Hausel's Betti-number
formula; hyperk\"ahler rotation and the topological Kempf--Ness theorem;
Brion's subgroup localization; constructible-sheaf formalism for vanishing
cycles (Maxim--Sch\"urmann, Massey); Kac polynomials and Crawley-Boevey's
lifting criterion; the polynomial-form CoHA and shuffle-to-difference-operator
maps (Jindal--Negu\c t, Tsymbaliuk).  The paper cites Krylov--Shlykov
\cite[Remark 1.15]{KS25} for related deformation arguments and
Dumanski--Krylov \cite[Prop.~2.3]{DK25} for the prequotient localization
viewpoint.

\subsection*{What appears to be new}
\begin{enumerate}[leftmargin=2em]
\item The \emph{stable-pair vanishing-cycle test} ($\S\S3$--$4$): a
finite-dimensional device turning a Coulomb relation into a statement
about sections of bundles on $\PP^1$, where stability of the value at
$\infty$ forces the degree inequality that kills the first step of every
neutral word.
\item The \emph{added-loop Laurent model} ($\S5$): one loop per vertex
cancels the Grassmannian Euler denominators, giving a polynomial model in
which an explicit multipartition filtration and cut ideals bound $B_0$.
\item The \emph{reachability stratification} ($\S6$) that computes the cut
quotients and matches Hausel's $q$-series exactly.
\item The \emph{homogeneity descent} ($\S7$) that removes the auxiliary
loop without any genericity assumption.
\item The \emph{complementary-monopole argument} ($\S8$) for the
cocharacter-versus-torus comparison, using Kac support plus regularity.
\end{enumerate}
\subsection*{Verification status}
After these notes were drafted, each of \S\S2--8 was read independently by
a separate Claude instance acting as an adversarial referee, recomputing
every formula it could and checking citations against the sources where
they could be fetched (BFN18, Weekes, Hoskins, McGerty--Nevins,
Tsymbaliuk, Jindal--Negu\c t).  \emph{No mathematical error and no
load-bearing gap was found.}  A dozen explicit computations (Jordan
$v=2$, $w=1,2$; $A_1$ and $A_2$ with loops; the $\mathrm{GL}_2$ loop
cancellation; the Tsymbaliuk kernel and sign table; the central splitting
in \S8) all agree.  The issues found are presentational: one convention
that should be stated (the grading is BFN's $N_{\mathcal O}$-relative one,
not the symmetric $2\Delta(\lambda)$; they differ by a function of charge),
two steps in \S4 that deserve a paragraph each (equivariant
Thom--Sebastiani; compatibility of two vanishing-cycle identifications,
which in any case only matters up to a unit), and a handful of citations
to make precise.  The full list is in \texttt{VERIFICATION.md} in the
repository.  This is a careful machine reading, not a human referee
report; it raises the prior that the paper is correct without settling
it.

\subsection*{Scope}
The theorem is for $G=\prod\GL(V_i)$ and $N$ of quiver type.  Loops,
multiple arrows, arbitrary flavor, and all regular stabilities are allowed.
Not covered: the quantized statement; other reductive gauge groups (the
paper cites Hoang--Krylov--Matvieievskyi, Ex.~6.7, for a failure outside
the quiver setting); $K$-theory.

\subsection*{A question}
The added-loop trick is a special case of adding a copy of the adjoint
representation: for any $G$, a massless adjoint summand contributes to the
abelianized monopole of \emph{any} coweight $\lambda$ the factor
$\prod_\alpha\alpha^{(-\langle\alpha,\lambda\rangle)_+}$, which cancels the Euler
class of $T_{t^\lambda}\Gr^\lambda_G$ up to sign.  So the Laurent-model step
and the mass-deformation step of $\S\S5,7$ look as if they could be run for
general $(G,N)$.  The other ingredients---minuscule generation, Kirwan
surjectivity, the cut-ideal count against Hausel's formula, and the
Kac-support argument of Part II---are all quiver-specific as written, and
the hypothesis of regularity rarely holds when $G$ has small center.
Whether some version of the argument extends is, as far as these notes go,
open.

\begin{thebibliography}{9}
\bibitem{OAI-Hikita} OpenAI, \emph{The equivariant cohomological Hikita
conjecture for arbitrary quivers}, OpenAI Math Release preprint, 24 September 2026.
\url{https://github.com/openai/math/blob/main/preprints/The-equivariant-cohomological-Hikita-conjecture-for-arbitrary-quivers-September-24-2026/main.pdf}
\bibitem{DK25} I.~Dumanski, V.~Krylov, \emph{$K$-theoretic Hikita conjecture for quiver gauge theories}.
\bibitem{Hausel06} T.~Hausel, \emph{Betti numbers of holomorphic symplectic quotients via arithmetic Fourier transform}, PNAS 2006.
\bibitem{Hikita17} T.~Hikita, \emph{An algebro-geometric realization of the cohomology ring of Hilbert scheme of points in the affine plane}, IMRN 2017.
\bibitem{KS25} V.~Krylov, P.~Shlykov, \emph{Hikita--Nakajima conjecture for the Gieseker variety}.
\bibitem{MN18} K.~McGerty, T.~Nevins, \emph{Kirwan surjectivity for quiver varieties}, Invent.\ Math.\ 2018.
\bibitem{Weekes19} A.~Weekes, \emph{Generators for Coulomb branches of quiver gauge theories}.
\end{thebibliography}

\end{document}
